\begin{tabular}{@{}>{\bfseries}p{0.12\textwidth}p{0.273\textwidth}p{0.273\textwidth}p{0.273\textwidth}@{}}
\toprule
 &\multicolumn{3}{c}{\textbf{Fermentation}}\\
 \cmidrule(rl){2-4}
 & \thead{Zu kurz} & \thead{Zu lang} & \thead{Perfekt} \\ \midrule
Krumen\-textur & Ungebackene, gummiartige Stellen am Boden des Brotes.
              & Krume kann als gummiartig empfunden werden, da das meiste Gluten abgebaut wurde.
              & Krume gleichmäßig gebacken. Krume kann als feucht, aber nicht gummiartig empfunden werden.
              \\
Alveolen       & Übermäßig große Alveolen in der Krume (,,Krater'').
              & Viele winzige, gleichmäßig verteilte Alveolen.
              & Alveolen gleichmäßig verteilt, keine ,,Krater''.
              \\
Geschmack         & Blasser, neutraler Geschmack.
              & Stark säuerliches Geschmacksprofil. Die Säure überwiegt beim Schmecken.
              & Ausgewogenes Geschmacksprofil, nicht zu mild, aber auch nicht zu sauer.
              Je nach Anstellgut essig- oder milchsäurebetonte Noten.
              \\
Textur       & Insgesamt schlechte Textur.
              & Gute Konsistenz, Krume ist nicht so locker wie möglich.
              & Großartige Kombination von Texturen.
              \\
Ofentrieb   & Vertikaler Ofentrieb, hauptsächlich durch verdampfendes Wasser, das den Teig aufbläht.
              & Sehr flache, pfannkuchenartige Struktur nach dem Backen.
              & Großartiger vertikaler Ofentrieb. Teig wächst mehr in die Höhe als in die Breite.
              \\ \bottomrule
\end{tabular}
